%% title: NLS Asymptotics - Derivative of the FOC and the Scaled Step
\documentclass[11pt]{article}
\usepackage{metricsstyle}

\begin{document}

\lessonheader{12}{Derivation supplement: the derivative of the first-order condition and the $\sqrt{n}$-scaled step}{}

\runin{Scope} These notes expand Section 3 of Lesson 12: how the derivative of the
NLS first-order condition is computed, and how the $\sqrt{n}$-scaled expansion
delivers asymptotic normality. Notation follows Lessons 10--12.

\section{Setup and notation}

\[ y = x(\beta) + u, \qquad u \sim (0,\sigma^{2}I_{n}),
   \qquad x(\beta) \ \text{an } n\times 1 \text{ vector}. \]
Write
\[ X(\beta) = \frac{\partial x(\beta)}{\partial\beta'} \ \ (n\times K), \qquad
   X_{t} = \frac{\partial x_{t}(\beta_{0})}{\partial\beta'} \ \ (1\times K), \qquad
   H_{t} = \frac{\partial^{2}x_{t}(\beta_{0})}{\partial\beta\,\partial\beta'}
           \ \ (K\times K). \]
Then $X_{0}=X(\beta_{0})$ stacks the rows $X_{t}$, and $u=y-x_{0}$.

\section{Where $\partial[X(\beta)'(y-x(\beta))]/\partial\beta$ comes from}

\topic{The function being differentiated}
Let $f(\beta)=X(\beta)'u(\beta)$ with $u(\beta)=y-x(\beta)$. Its $j$-th entry is
$f_{j}(\beta)=\sum_{t}X_{tj}(\beta)\,u_{t}(\beta)$. Differentiating in
$\beta_{l}$ with the product rule,
\[ \frac{\partial f_{j}}{\partial\beta_{l}}
   = \underbrace{\sum_{t}\frac{\partial X_{tj}}{\partial\beta_{l}}\,u_{t}}_{(A)}
   + \underbrace{\sum_{t}X_{tj}\,\frac{\partial u_{t}}{\partial\beta_{l}}}_{(B)}. \]

\topic{Term (B): $u$ moves, $X$ fixed}
Since $y$ does not depend on $\beta$, $\partial u_{t}/\partial\beta_{l}=-X_{tl}$,
so
\[ (B) = -\sum_{t}X_{tj}X_{tl} = -(X'X)_{jl}. \]

\topic{Term (A): $X$ moves, $u$ fixed}
Since $X_{tj}=\partial x_{t}/\partial\beta_{j}$,
\[ (A) = \sum_{t}u_{t}\,\frac{\partial^{2}x_{t}}{\partial\beta_{j}\,\partial\beta_{l}}
       = \sum_{t}u_{t}\,(H_{t})_{jl}. \]
This is what the lecture writes as
$\frac{\partial X(\beta)'}{\partial\beta}\big|_{\beta_{0}}(y-x_{0})$: the
residual-weighted Hessian of $x(\beta)$.

\topic{Evaluating at $\beta_{0}$}
At $\beta_{0}$ we have $u=y-x_{0}$, so
\[ \frac{\partial}{\partial\beta}
   \Bigl[X(\beta)'\bigl(y-x(\beta)\bigr)\Bigr]_{\beta_{0}}
   = \sum_{t=1}^{n}u_{t}H_{t} - X_{0}'X_{0}. \]
In the linear model $x(\beta)=X\beta$, so $H_{t}=0$ and only $-X'X$ remains.

\section{The $\sqrt{n}$-scaled equation, written out}

Taylor-expanding the first-order condition $X(\hat\beta)'(y-x(\hat\beta))=0$
around $\beta_{0}$ and multiplying by $1/\sqrt{n}$ gives
\[ 0 = a_{n} + (A_{n}-S_{n})\,z_{n} + o_{p}(1), \]
with the four objects
\begin{align*}
 a_{n} &= \frac{X_{0}'u}{\sqrt{n}}
        = \frac{1}{\sqrt{n}}\sum_{t=1}^{n}X_{t}'u_{t}
 && (K\times1), \qquad (a_{n})_{j}
    = \frac{1}{\sqrt{n}}\sum_{t}\frac{\partial x_{t}}{\partial\beta_{j}}\,u_{t},\\[3pt]
 S_{n} &= \frac{X_{0}'X_{0}}{n}
        = \frac{1}{n}\sum_{t=1}^{n}X_{t}'X_{t}
 && (K\times K), \qquad (S_{n})_{jl}
    = \frac{1}{n}\sum_{t}\frac{\partial x_{t}}{\partial\beta_{j}}
      \frac{\partial x_{t}}{\partial\beta_{l}},\\[3pt]
 A_{n} &= \frac{1}{n}\sum_{t=1}^{n}u_{t}H_{t}
 && (K\times K), \qquad (A_{n})_{jl}
    = \frac{1}{n}\sum_{t}u_{t}\,
      \frac{\partial^{2}x_{t}}{\partial\beta_{j}\,\partial\beta_{l}},\\[3pt]
 z_{n} &= \sqrt{n}\,(\hat\beta-\beta_{0})
 && (K\times1).
\end{align*}

\section{Limits of each object}

\topic{$S_{n}\xrightarrow{\ p\ }\Sigma_{XX}$}
This is the identification assumption.

\topic{$A_{n}\xrightarrow{\ p\ }0$}
The lecture says ``its expectation is zero'', which is not sufficient by itself.
Treat $H_{t}$ as nonrandom (or condition on it). Then
\[ \E\bigl[(A_{n})_{jl}\bigr]
   = \frac{1}{n}\sum_{t}\E(u_{t})\,(H_{t})_{jl} = 0, \]
\[ \Var\bigl[(A_{n})_{jl}\bigr]
   = \frac{\sigma^{2}}{n^{2}}\sum_{t}(H_{t})_{jl}^{2}
   = \frac{\sigma^{2}}{n}\cdot
     \underbrace{\frac{1}{n}\sum_{t}(H_{t})_{jl}^{2}}_{=\,O(1)}
   = O(1/n). \]
Mean zero and variance tending to zero give $(A_{n})_{jl}\xrightarrow{p}0$ by
Chebyshev, for every $(j,l)$. The $O(1)$ bound on
$\frac{1}{n}\sum_{t}(H_{t})_{jl}^{2}$ is a regularity condition that must be
assumed.

\topic{$a_{n}=O_{p}(1)$}
\[ \E(a_{n}) = 0, \qquad
   \Var(a_{n}) = \frac{1}{n}\sum_{t}X_{t}'\sigma^{2}X_{t}
   = \sigma^{2}S_{n} \longrightarrow \sigma^{2}\Sigma_{XX} < \infty. \]
A bounded variance gives $a_{n}=O_{p}(1)$ by Chebyshev.

\topic{$z_{n}=O_{p}(1)$}
Consistency of $\hat\beta$ alone does not give this: $\hat\beta-\beta_{0}
\xrightarrow{p}0$ says nothing about the $\sqrt{n}$ rate. Rearrange the scaled
equation:
\[ z_{n} = (S_{n}-A_{n})^{-1}\bigl(a_{n}+o_{p}(1)\bigr). \]
Since $S_{n}-A_{n}\xrightarrow{p}\Sigma_{XX}$ is invertible, the continuous
mapping theorem gives $(S_{n}-A_{n})^{-1}\xrightarrow{p}\Sigma_{XX}^{-1}=O(1)$.
Multiplying by $a_{n}=O_{p}(1)$ gives $z_{n}=O_{p}(1)$.

\section{Why $\dfrac{X_{0}'u}{\sqrt{n}}-\Sigma_{XX}\sqrt{n}(\hat\beta-\beta_{0})
\xrightarrow{p}0$}

From the scaled equation, isolate $a_{n}-S_{n}z_{n}$:
\[ a_{n}-S_{n}z_{n} = -A_{n}z_{n} - o_{p}(1). \]
Add and subtract $\Sigma_{XX}z_{n}$, so that the left side becomes
$a_{n}-\Sigma_{XX}z_{n}$ plus the gap $(\Sigma_{XX}-S_{n})z_{n}$:
\[ a_{n}-\Sigma_{XX}z_{n}
   = -\underbrace{A_{n}}_{\xrightarrow{\ p\ }0}z_{n}
   + \underbrace{(S_{n}-\Sigma_{XX})}_{\xrightarrow{\ p\ }0}z_{n}
   - o_{p}(1). \]

\topic{Entry by entry}
Row $j$ of this equation is
\[ (a_{n})_{j}-\sum_{l=1}^{K}(\Sigma_{XX})_{jl}\,z_{n,l}
   = -\sum_{l=1}^{K}(A_{n})_{jl}\,z_{n,l}
   + \sum_{l=1}^{K}\bigl[(S_{n})_{jl}-(\Sigma_{XX})_{jl}\bigr]z_{n,l}
   - o_{p}(1). \]
Every summand on the right is
\[ \underbrace{(\text{scalar}\xrightarrow{p}0)}_{o_{p}(1)}
   \times \underbrace{z_{n,l}}_{O_{p}(1)} = o_{p}(1), \]
and a finite sum ($K$ terms) of $o_{p}(1)$ terms is $o_{p}(1)$. This holds for
every $j=1,\dots,K$, so the whole vector converges:
\[ \boxed{\ \frac{X_{0}'u}{\sqrt{n}}
   - \Sigma_{XX}\sqrt{n}(\hat\beta-\beta_{0})
   \xrightarrow{\ p\ } 0\ } \]

\section{Where this leads}

Since $X_{0}'u/\sqrt{n}\xrightarrow{d}N(0,\sigma^{2}\Sigma_{XX})$ by the CLT,
Slutsky gives $\Sigma_{XX}\sqrt{n}(\hat\beta-\beta_{0})
\xrightarrow{d}N(0,\sigma^{2}\Sigma_{XX})$, hence
\[ \sqrt{n}(\hat\beta-\beta_{0})
   \xrightarrow{\ d\ } N\bigl(0,\ \sigma^{2}\Sigma_{XX}^{-1}
   \Sigma_{XX}\Sigma_{XX}^{-1}\bigr)
   = N\bigl(0,\ \sigma^{2}\Sigma_{XX}^{-1}\bigr). \]

\begin{transcriptnotes}
 \item This is a derivation supplement to Lesson 12, not a transcription of the
       lecture. The notation, title block, numbered sections and topic headings
       follow Lessons 10--12 and the shared \texttt{metricsstyle} package.
 \item The lecture's derivative of the first-order condition,
       $\partial X'/\partial\beta\,u$, is ambiguous as written; it is computed
       here as $\sum_{t}u_{t}H_{t}(\beta_{0})$. The mean-zero condition and the
       Chebyshev argument for $A_{n}\xrightarrow{p}0$ are made explicit, since
       $\E(u_{t}H_{t})=0$ alone does not give the limit.
 \item The $\sqrt{n}$-scaled equation, the $O_{p}(1)$ arguments and the
       Slutsky step expand Section 3 of Lesson 12; the boxed result is the
       ingredient used there to obtain
       $\sqrt{n}(\hat\beta-\beta_{0})\xrightarrow{d}N(0,\sigma^{2}
       \Sigma_{XX}^{-1})$.
\end{transcriptnotes}

\end{document}
